\MathBookSourcePage{240}
\subsection{$L^p$ 估计(正则系数)}
\label{sec:ch08-lp-estimates-regular-coefficients}
\setcounter{thmcount}{0}
\setcounter{equation}{0}

当 $2\leq p\leq\infty$ 时, 将 $p$ 替换定理~\ref{thm:ch08-main-max-norm-stability} 中的 $\infty$ 的证明见习题~\ref{exr:ch08-01}. 该估计可由下面的对偶论证推广到 $1<p\leq2$. 令 $q\in C_0^\infty(\overline\Omega)$, 且 $1/p+1/q=1$. 由 Simader (1972),
\[
\MathBookFormulaID{formula-ch08-lp-inf-sup-estimate}
\frac1C\|u_h\|_{W_p^1(\Omega)}
\leq\sup_{0\neq v\in W_q^1(\Omega)}\frac{a(u_h,v)}{\|v\|_{W_q^1(\Omega)}}
+\|u_h\|_{L^p(\Omega)}.
\]
若 $v_h$ 表示 $v$ 关于 $a(\cdot,\cdot)$ 的投影, 则应用习题~\ref{exr:ch08-01} 和 Hölder 不等式得到
\[
\MathBookFormulaID{formula-ch08-lp-discrete-stability-preliminary}
\|u_h\|_{W_p^1(\Omega)}\leq C\bigl(\|u\|_{W_p^1(\Omega)}+\|u_h\|_{L^p(\Omega)}\bigr).
\]

对 $\|u_h\|_{L^p(\Omega)}$ 的界可由对偶性得到, 只需略微修改定理~\ref{thm:ch05-poisson-l2-error} 的证明. 考虑变分问题: 求 $w\in V$ 使得
\MBSourceItem{1}
\begin{equation}
\MathBookFormulaID{formula-ch08-dual-lp-problem}
\label{eq:ch08-dual-lp-problem}
a(v,w)=(f,v)
\qquad\forall v\in V.
\end{equation}
当 $1<q\leq\infty$ 时, 由式~\eqref{eq:ch08-primal-regularity} 有
\MBSourceItem{2}
\begin{equation}
\MathBookFormulaID{formula-ch08-dual-lp-regularity}
\label{eq:ch08-dual-lp-regularity}
\|w\|_{W_q^1(\Omega)}\leq C\|f\|_{L^q(\Omega)},
\end{equation}
其中需有 $\mu>d$(参见习题~\ref{exr:ch08-14}). 由此和 Sobolev、Hölder 不等式, 取 $f=|u_h|^{p-1}\operatorname{sign}(u_h)$, 可得
\[
\MathBookFormulaID{formula-ch08-lp-discrete-lp-bound}
\|u_h\|_{L^p(\Omega)}\leq C\|u\|_{W_p^1(\Omega)}
\qquad\forall 1<p\leq\infty.
\]
合并上述估计得到下述结论.

\MBSourceItem{3}
\begin{Theorem}
\label{thm:ch08-lp-stability-regular-coefficients}
在定理~\ref{thm:ch08-main-max-norm-stability} 的条件下, 并假设 $a_{ij}\in C^1(\overline\Omega)$, 则存在 $h_0>0$ 和 $C<\infty$, 使得对 $0<h<h_0$,
\[
\MathBookFormulaID{formula-ch08-lp-stability-regular-coefficients}
\|u_h\|_{W_p^1(\Omega)}\leq C\|u\|_{W_p^1(\Omega)}
\qquad\forall 1<p\leq\infty.
\]
\end{Theorem}

由上一结果应用于 $u-v$, 再用三角不等式, 当然可得
\MBSourceItem{4}
\begin{equation}
\MathBookFormulaID{formula-ch08-lp-best-approximation}
\label{eq:ch08-lp-best-approximation}
\|u-u_h\|_{W_p^1(\Omega)}\leq C\inf_{v\in V_h}\|u-v\|_{W_p^1(\Omega)}
\qquad\forall1<p\leq\infty.
\end{equation}
把上述结果与对偶技巧结合, 还可推出(证明留作习题~\ref{exr:ch08-13})
\MBSourceItem{5}
\begin{equation}
\MathBookFormulaID{formula-ch08-lp-lower-order-error}
\label{eq:ch08-lp-lower-order-error}
\|u-u_h\|_{L^p(\Omega)}\leq Ch\|u-u_h\|_{W_p^1(\Omega)}
\qquad\forall\mu'<p<\infty,
\end{equation}
其中 $1/\mu+1/\mu'=1$(参见式~\eqref{eq:ch08-primal-regularity}). $p=\infty$ 的情形因式~\eqref{eq:ch08-primal-regularity} 对 $p=1$ 缺少正则性而较特殊. 分片线性元的最佳阶逼近至少在二维中具有轻微的对数损失(Scott 1976):
\[
\MathBookFormulaID{formula-ch08-linfty-logarithmic-error}
\|u-u_h\|_{L^\infty(\Omega)}\leq Ch|\log h|\|u-u_h\|_{W_\infty^1(\Omega)}.
\]
对高次逼近可证明(Scott 1976)
\[
\MathBookFormulaID{formula-ch08-higher-order-linfty-error}
\|u-u_h\|_{L^\infty(\Omega)}\leq Ch\|u-u_h\|_{W_\infty^1(\Omega)}.
\]

这里不再推导正则系数情形的更多 $L^p$ 估计, 下一节转而说明即使系数很粗糙也能得到某些估计.
